% Section: Glossary C.6: connections and curvature | Part: backmatter | Depends on: sec:conn, sec:lc, sec:curv, sec:synth3 | Status: done | Original lines: none (added by rewriter)
%!TEX root = ../../../main.tex
\subsubsection*{Connections and curvature}

Conventions as in Part~\ref{part:geometry}: $\nabla_{\partial_i}\partial_j=\Gamma^k{}_{ij}\partial_k$, $R(X,Y)Z=\nabla_X\nabla_YZ-\nabla_Y\nabla_XZ-\nabla_{[X,Y]}Z$,
$R(\partial_i,\partial_j)\partial_k=R^l{}_{ijk}\partial_l$ and $\mathbb R(X,Y,Z,W)=g(R(X,Y)Z,W)$.

% PROOF ADDED BY REWRITER: eq:christtransf is stated without derivation.
\begin{identity}[transformation of the Christoffel symbols; \eqref{eq:christtransf}]\label{id:gammatrans}
Under a change of chart $x\mapsto\tilde x$ the symbols transform by \eqref{eq:christtransf}.
\end{identity}
\begin{proof}
\Step{1}{new coordinate fields} By the chain rule, $\tilde\partial_{\tilde i}=\frac{\partial x^i}{\partial\tilde x^{\tilde i}}\partial_i$.
\Step{2}{expand with the axioms} By $C^\infty$-linearity in the direction and the Leibniz rule,
\[
\nabla_{\tilde\partial_{\tilde i}}\tilde\partial_{\tilde j}=\frac{\partial x^i}{\partial\tilde x^{\tilde i}}\Bigl(\partial_i\Bigl(\frac{\partial x^k}{\partial\tilde x^{\tilde j}}\Bigr)\partial_k+\frac{\partial x^j}{\partial\tilde x^{\tilde j}}\Gamma^k{}_{ij}\partial_k\Bigr).
\]
\Step{3}{second derivatives} $\frac{\partial x^i}{\partial\tilde x^{\tilde i}}\partial_i\bigl(\frac{\partial x^k}{\partial\tilde x^{\tilde j}}\bigr)=\frac{\partial^2x^k}{\partial\tilde x^{\tilde i}\partial\tilde x^{\tilde j}}$ (chain rule again).
\Step{4}{back to the new basis} Write $\partial_k=\frac{\partial\tilde x^{\tilde k}}{\partial x^k}\tilde\partial_{\tilde k}$ and read off the coefficient of $\tilde\partial_{\tilde k}$.
\end{proof}

% PROOF ADDED BY REWRITER: the text says "derive it by expanding".
\begin{identity}[Riemann tensor in coordinates; \eqref{eq:riemanncomp}]\label{id:riemanncomp}
$R^l{}_{ijk}=\partial_i\Gamma^l{}_{jk}-\partial_j\Gamma^l{}_{ik}+\Gamma^s{}_{jk}\Gamma^l{}_{is}-\Gamma^s{}_{ik}\Gamma^l{}_{js}$.
\end{identity}
\begin{proof}
\Step{1}{one term} $\nabla_{\partial_i}(\nabla_{\partial_j}\partial_k)=\nabla_{\partial_i}(\Gamma^s{}_{jk}\partial_s)=\partial_i\Gamma^l{}_{jk}\,\partial_l+\Gamma^s{}_{jk}\Gamma^l{}_{is}\,\partial_l$ (Leibniz rule; the dummy index renamed).
\Step{2}{antisymmetrise} Subtract the same expression with $i\leftrightarrow j$; the term $\nabla_{[\partial_i,\partial_j]}$ is absent because coordinate fields commute.
\end{proof}

% PROOF ADDED BY REWRITER: the text says "It follows from the Jacobi identity".
\begin{identity}[first Bianchi identity, torsion-free case]\label{id:bianchi1}
If $\nabla$ is torsion-free, $R(X,Y)Z+R(Y,Z)X+R(Z,X)Y=0$; also $R(X,Y)=-R(Y,X)$ for every connection.
\end{identity}
\begin{proof}
\Step{1}{antisymmetry} The definition of $R(X,Y)Z$ changes sign under $X\leftrightarrow Y$, since $[Y,X]=-[X,Y]$.
\Step{2}{reduce to coordinate fields} $R$ is a tensor, so it suffices to take $X,Y,Z$ coordinate fields. Then all brackets vanish, and torsion-freeness reads $\nabla_XY=\nabla_YX$.
\Step{3}{regroup} The cyclic sum is
$\nabla_X(\nabla_YZ-\nabla_ZY)+\nabla_Y(\nabla_ZX-\nabla_XZ)+\nabla_Z(\nabla_XY-\nabla_YX)$, and every bracket vanishes by Step~2.
\end{proof}

% PROOF ADDED BY REWRITER: the metric-case properties (a), (b) are stated without proof.
\begin{identity}[symmetries of $\mathbb R$ for Levi-Civita]\label{id:riemannsym}
$\mathbb R(X,Y,Z,W)=-\mathbb R(X,Y,W,Z)$ and $\mathbb R(X,Y,Z,W)=\mathbb R(Z,W,X,Y)$.
\end{identity}
\begin{proof}
\Step{1}{it suffices that $\mathbb R(X,Y,Z,Z)=0$} Expanding $\mathbb R(X,Y,Z+W,Z+W)=0$ gives the antisymmetry in the last pair.
\Step{2}{use $\nabla g=0$} Take $X,Y$ coordinate fields, so $[X,Y]=0$. Then
$g(\nabla_X\nabla_YZ,Z)=X\,g(\nabla_YZ,Z)-g(\nabla_YZ,\nabla_XZ)=\tfrac12XY\,g(Z,Z)-g(\nabla_YZ,\nabla_XZ)$.
\Step{3}{antisymmetrise in $X,Y$} The last term is symmetric in $X,Y$, and $XY-YX=[X,Y]=0$ on functions. So $\mathbb R(X,Y,Z,Z)=0$.
\Step{4}{pair symmetry} Write the first Bianchi identity (Identity~\ref{id:bianchi1}, paired with $W$) for the cyclic orders
$(X,Y,Z;W)$, $(Y,Z,W;X)$, $(Z,W,X;Y)$, $(W,X,Y;Z)$, and add the four. By the two antisymmetries, eight of the twelve terms cancel in pairs, leaving
$-2\,\mathbb R(X,Z,Y,W)+2\,\mathbb R(Y,W,X,Z)=0$. This is the pair symmetry, with $(X,Z)$ and $(Y,W)$ as the pairs.
\end{proof}

\begin{identity}[$\dd$ and $\Lie_X$ through $\nabla$; Proposition~\ref{prop:threederivs}]\label{id:threederivs}
For Levi-Civita: $(\dd\omega)(X,Y)=(\nabla_X\omega)(Y)-(\nabla_Y\omega)(X)$ and $(\Lie_Xg)_{ij}=\nabla_iX_j+\nabla_jX_i$.
\end{identity}
\begin{proof}
\Step{1}{(a)} $(\nabla_X\omega)(Y)=X(\omega(Y))-\omega(\nabla_XY)$. Antisymmetrise and use $\nabla_XY-\nabla_YX=[X,Y]$; Identity~\ref{id:dinvariant} gives $\dd\omega(X,Y)$.
\Step{2}{(b)} $\Lie_Xg(Y,Z)=X(g(Y,Z))-g([X,Y],Z)-g(Y,[X,Z])$. Expand $X(g(Y,Z))$ with $\nabla g=0$ and the brackets with torsion-freeness:
$\Lie_Xg(Y,Z)=g(\nabla_YX,Z)+g(Y,\nabla_ZX)$, which in components is $\nabla_iX_j+\nabla_jX_i$.
\end{proof}

\begin{identity}[the divergences agree; Proposition~\ref{prop:divagree}]\label{id:divagree}
$\Gamma^i{}_{ij}=\partial_j\ln\sqrt{|g|}$, hence $\nabla_iX^i=\frac1{\sqrt{|g|}}\partial_i(\sqrt{|g|}X^i)=\operatorname{div}_{\vol_g}X$.
\end{identity}
\begin{proof}
\Step{1}{contract \eqref{eq:christmetric}} $\Gamma^i{}_{ij}=\frac12g^{il}(\partial_ig_{jl}+\partial_jg_{il}-\partial_lg_{ij})$; the first and last terms cancel after $i\leftrightarrow l$, leaving $\frac12g^{il}\partial_jg_{il}$.
\Step{2}{Jacobi's formula} $\det(g+\varepsilon h)=\det g\,\det(1+\varepsilon g^{-1}h)=\det g\,(1+\varepsilon\operatorname{tr}(g^{-1}h)+O(\varepsilon^2))$, so $\partial_j\ln|\det g|=g^{il}\partial_jg_{il}$.
\Step{3}{conclude} $\nabla_iX^i=\partial_iX^i+X^j\partial_j\ln\sqrt{|g|}=\frac1{\sqrt{|g|}}\partial_i(\sqrt{|g|}X^i)$, which is Identity~\ref{id:divmu}.
\end{proof}
